\chapter{Conclusão}\label{cap:conclusao}

Este trabalho desenvolveu um gerador de formas de onda por síntese digital direta, com núcleo em \ac{vhdl} no \ac{fpga} da placa DE2-115, controle pelo computador através do Virtual \ac{jtag}, placa analógica própria e interface gráfica em Python. O gerador produz senoides, rampas, funções sinc e formas arbitrárias de 0 a 262\,143~Hz, com resolução de 2,33~mHz, e dispensa qualquer hardware de comunicação além do cabo de gravação da placa.

O projeto digital foi organizado em blocos encapsulados, com uma entidade por função e um pacote único de constantes. Essa organização permitiu verificar cada bloco isoladamente: os 16 \eng{testbenches} auto-verificáveis, que vão da conversão de frequência à entidade do topo, terminaram sem erros. A separação entre o \ac{ip} do Virtual \ac{jtag} e a lógica de controle permitiu simular o caminho completo, do comando do computador até a amostra no \ac{dac}, que de outra forma só poderia ser verificado na placa. A conversão de frequência por multiplicação em ponto fixo, sem divisor, manteve o erro abaixo de 1~\ac{lsb} da palavra de sintonia em toda a faixa, e a travessia dos comandos entre os domínios de clock, feita com um único bit sincronizado, entregou todos os comandos testados sem perda nem repetição.

O projeto ocupa menos de 1\% dos elementos lógicos do \ac{fpga} e opera com folga de tempo de dez vezes no clock de 10~MHz. A simulação numérica do espectro confirmou que, com tabela de 1024 posições, o maior espúrio de truncamento fica em cerca de \(-60\)~dBc. Assim, o desempenho espectral é limitado principalmente pela resolução de 8~bits do DAC0800.

Na simulação da placa analógica, a saída reproduziu o seno gerado pelo \ac{vhdl} com a amplitude, o ganho e o atraso calculados. A distorção harmônica total foi de \(-70{,}9\)~dBc, e a maior componente indesejada, em \(-67\)~dBc após o filtro, é a imagem associada à resolução da tabela, de acordo com o limite de \(6{,}02\,P\). O estágio de saída classe AB, previsto no projeto da placa, foi retirado por distorção, e a saída passou a ser tomada do seguidor de tensão.

\pendente{Completar com os resultados dos ensaios em bancada com o código final.}

Como trabalhos futuros, propõem-se:
\begin{itemize}
	\item realizar os ensaios em bancada com o código final, medindo a frequência ao longo da faixa e o espectro da saída, para comparar o maior espúrio com o previsto;
	\item comparar a forma de onda e o espectro medidos na placa com os da simulação SPICE;
	\item projetar um estágio de saída com maior capacidade de corrente e baixa distorção, em substituição ao classe AB retirado;
	\item substituir o DAC0800 por um conversor de maior resolução e menor tempo de acomodação, aproveitando a margem de tempo do \ac{fpga} para elevar o clock e a frequência máxima;
	\item usar um filtro de reconstrução de ordem mais alta, para atenuar mais as imagens com a mesma banda passante;
	\item explorar a simetria do seno para reduzir a tabela, liberando memória para tabelas maiores ou para mais formas de onda;
	\item acrescentar modulações (amplitude, frequência e varredura) e um segundo canal com defasagem controlada.
\end{itemize}
